\documentclass[11pt,a4paper]{article}
\usepackage[margin=25mm,headheight=15pt]{geometry}
\usepackage[T1]{fontenc}
\usepackage[utf8]{inputenc}
\usepackage{lmodern,amsmath,amssymb,amsthm,mathtools,bm}
\usepackage{xcolor,graphicx,booktabs,tabularx,array,longtable}
\usepackage{microtype,enumitem,fancyhdr}
\usepackage{xurl}
\usepackage{hyperref}
\definecolor{accent}{HTML}{163B55}
\hypersetup{colorlinks=true,linkcolor=accent,citecolor=accent,urlcolor=accent,
 pdftitle={Invariant subspaces after complexification of a real irreducible operator},pdfauthor={Matthew J. Colbrook},
 pdfsubject={AIM 920: mathematical result and proof}}
\urlstyle{same}
\setlength{\parindent}{0pt}
\setlength{\parskip}{0.35em}
\setlength{\emergencystretch}{3em}
\setlist{itemsep=0.25em,topsep=0.3em}
\newtheorem{theorem}{Theorem}[section]
\newtheorem{lemma}[theorem]{Lemma}
\newtheorem{proposition}[theorem]{Proposition}
\newtheorem{corollary}[theorem]{Corollary}
\theoremstyle{remark}
\newtheorem{remark}[theorem]{Remark}
\newtheorem{definition}[theorem]{Definition}
\DeclareMathOperator{\Tr}{Tr}
\DeclareMathOperator{\supp}{supp}
\DeclareMathOperator{\Lip}{Lip}
\DeclareMathOperator{\Rea}{Re}
\newcommand{\R}{\mathbb R}
\newcommand{\C}{\mathbb C}
\newcommand{\E}{\mathbb E}
\newcommand{\N}{\mathbb N}
\newcommand{\ip}[2]{\langle #1,#2\rangle}
\pagestyle{fancy}
\fancyhf{}
\fancyhead[L]{\small\sffamily AIM 920}
\fancyhead[R]{\small\sffamily Matthew J. Colbrook}
\fancyfoot[C]{\thepage}
\renewcommand{\headrulewidth}{0.3pt}
\allowdisplaybreaks[2]
\title{\textbf{Invariant subspaces after complexification of a real irreducible operator}}
\author{Matthew J. Colbrook\\[3pt]
{\small Department of Applied Mathematics and Theoretical Physics}\\
{\small University of Cambridge, United Kingdom}\\
{\small\href{mailto:mjc249@maths.cam.ac.uk}{\texttt{mjc249@maths.cam.ac.uk}}}}
\date{3 November 2026}
\begin{document}
\maketitle
\begin{abstract}
We construct a real irreducible bounded operator whose complexification has exactly two proper nonzero closed invariant subspaces. A real polynomial approximation recovers the complex structure, and spectral separation determines the complexified invariant-subspace lattice. We compare the construction with the original conjecture and a later proposition.
\end{abstract}
\noindent\textbf{Repository target.}
This manuscript concerns AIM problem 920 at revision
\texttt{5d7e26e8444e}. The precise scope is stated below.
\medskip

\section{Main result and proof}
\label{sec:920}

The catalogue asks whether a bounded operator on an infinite-dimensional
separable real Banach space, having no proper nonzero closed real
invariant subspace, must retain that property after complexification.
This is Conjecture~3 in \cite{AAST2005}.  The following construction
gives a negative answer.  It uses the classical existence of an operator
on a complex Banach space with no proper nonzero closed complex invariant
subspace; the passage to the required real operator is elementary.

\begin{theorem}\label{thm:920}
There are an infinite-dimensional separable real Banach space $X$ and a
bounded real-linear operator $T$ such that
\[
 \operatorname{Lat}_{\mathbb R}(T)=\{\{0\},X\},
 \qquad
 \operatorname{Lat}_{\mathbb C}(T_{\mathbb C})
 =\{\{0\},E_+,E_-,X_{\mathbb C}\},
\]
where $E_+,E_-$ are nonzero closed complementary complex subspaces.
One may take $X$ to be the real space underlying $\ell_1(\mathbb C)$,
or equivalently transfer the example to $\ell_1(\mathbb R)$.
The operator can also be chosen invertible and of norm one.
\end{theorem}

\begin{lemma}\label{lem:920-recovery}
Let $Y$ be a complex Banach space, let $X=Y_{\mathbb R}$ be its
underlying real space, and let $Jx=ix$ denote its original complex
structure.  If $S\in\mathcal L_{\mathbb C}(Y)$ and $\|S\|\leq1/10$,
then $J$ lies in the norm-closed unital real algebra generated by
$T=J+S$.
\end{lemma}
\begin{proof}
We have $J^2=-I$, $\|J\|=1$, and $JS=SJ$.  Define
\[
 Q=I-JS,\qquad E=-T^2-I=-2JS-S^2.
\]
Then $T=JQ$, $Q^2=I+E=-T^2$, $\|Q-I\|\leq1/10$, and
$\|E\|\leq21/100<1$.  In the norm-closed real algebra generated by
$T$, the real binomial series
\[
 U=\sum_{n=0}^{\infty}\binom{1/2}{n}E^n
\]
converges absolutely and satisfies $U^2=I+E=Q^2$.  Moreover
\[
 \|U-I\|\leq
 \sum_{n\geq1}\left|\binom{1/2}{n}\right|(21/100)^n
 =1-\sqrt{79/100}<\frac{12}{100}.
\]
The operators $U,Q$ commute, because $Q$ commutes with $E$.  Also
$\|(U+Q)-2I\|<22/100<2$, so $U+Q$ is invertible.  Factoring
$U^2-Q^2$ yields $U=Q$.  Since $\|U-I\|<1$, the inverse $U^{-1}$
is a Neumann series in the same closed real algebra.  Therefore
$J=TU^{-1}$ belongs to that algebra, as asserted.
\end{proof}

\begin{proof}[Proof of Theorem~\ref{thm:920}]
Choose a bounded complex-linear operator $R$ on $Y=\ell_1(\mathbb C)$
with no proper nonzero closed complex invariant subspace.  This is the
classical Read construction \cite{Read1986}.  For an explicit modern
primary source covering the complex field, the introduction of
\cite{GR2014} allows real or complex spaces, and its Theorem~1.1 gives
the stronger absence of nontrivial invariant closed subsets on
$\ell_1(\mathbb C)$: take $Z=\ell_1(\mathbb C)$ and use the isometry
$\ell_1(Z)\cong\ell_1(\mathbb C)$.

Multiply $R$ by a nonzero real scalar to obtain $S$ with
$\|S\|\leq1/10$.  Its complex invariant subspaces are unchanged.
Regard $Y$ as the real space $X$ and set $T=J+S$.
If a closed real subspace $M$ is $T$-invariant, it is invariant under
every real polynomial in $T$, and therefore under their operator-norm
limits.  Lemma~\ref{lem:920-recovery} gives $J(M)\subseteq M$.
Thus $M$ is a complex subspace of $Y$, and
$S(M)=(T-J)(M)\subseteq M$.  By the choice of $S$, either $M=\{0\}$
or $M=X$.  This proves real irreducibility.

Now form the external complexification $X_{\mathbb C}=X+iX$, with
the Taylor norm specified in the catalogue.  The external scalar $i$
is distinct from the original complex structure $J$ on $X$.
The bounded complex-linear extension $J_{\mathbb C}$ commutes with
$T_{\mathbb C}$ and satisfies $J_{\mathbb C}^2=-I$.  Hence
\[
 P_+=\frac{I-iJ_{\mathbb C}}2,\qquad
 P_-=\frac{I+iJ_{\mathbb C}}2
\]
are complementary bounded projections commuting with $T_{\mathbb C}$.
Their ranges are
\[
 E_+=\ker(J_{\mathbb C}-iI)
       =\{x-iJx:x\in X\},\qquad
 E_-=\ker(J_{\mathbb C}+iI)
       =\{x+iJx:x\in X\}.
\]
For every $x\ne0$ both displayed vectors are nonzero.  Thus these are
proper nonzero closed complex invariant subspaces, and
$X_{\mathbb C}=E_+\oplus E_-$.

We also identify the entire invariant-subspace lattice.
Complexification is norm-continuous for the Taylor norm; indeed
$\|A_{\mathbb C}\|=\|A\|$ for every bounded real operator $A$.
Lemma~\ref{lem:920-recovery} therefore implies that
$J_{\mathbb C}$ is a norm limit of real polynomials in $T_{\mathbb C}$.
Every closed complex $T_{\mathbb C}$-invariant subspace $N$ is consequently
invariant under $P_+,P_-$, and
\[
 N=(N\cap E_+)\oplus(N\cap E_-).
\]
The map $x\mapsto x-iJx$ is a complex-linear topological isomorphism
from $Y$ to $E_+$, intertwining $S+iI$ with
$T_{\mathbb C}|_{E_+}$.  The analogous map to $E_-$ is conjugate-linear,
so its restriction is likewise complex irreducible.  Each intersection
is therefore either zero or the whole corresponding $E_\pm$, proving
the asserted four-element lattice.

Finally $T=JQ$ is invertible because $\|Q-I\|\leq1/10$.
Multiplying $T$ by $\|T\|^{-1}$ changes neither invariant-subspace
lattice.  The real-linear identification
\[
 (a_1,b_1,a_2,b_2,\ldots)\longmapsto(a_1+ib_1,a_2+ib_2,\ldots)
\]
is an isomorphism $\ell_1(\mathbb R)\to Y_{\mathbb R}$, since
\[
 \sum_n\sqrt{a_n^2+b_n^2}
 \leq\sum_n(|a_n|+|b_n|)
 \leq\sqrt2\sum_n\sqrt{a_n^2+b_n^2}.
\]
Transporting $T$ through this isomorphism gives the claimed example
on the classical real $\ell_1$ space.
\end{proof}

\subsection*{An independent consequence of an existing theorem}

There is also a shorter route using the stronger result already proved
in \cite[Theorem~1.1]{GR2014}.  On $Y=\ell_1(\mathbb C)$ it supplies
an operator $R$ for which every nonzero vector has dense forward orbit.
The topology is unchanged on passing to the underlying real space.
Therefore every nonzero closed real $R$-invariant subspace contains a
dense orbit and equals $Y_{\mathbb R}$.  Nevertheless the external
complexification of $R$ preserves the nonzero proper eigenspaces of
the original complex structure $J_{\mathbb C}$, exactly as above.
Thus a negative answer to the stated catalogue problem is also a
direct consequence of that existing theorem.  This observation is
not a claim that the existence theorem or this consequence is new.
The preceding proof needs only the weaker classical invariant-subspace
counterexample and additionally identifies the complete complexified
invariant-subspace lattice.

\subsection*{Conflict with a cited proposition}

There is a specific literature conflict that should accompany this
answer.  Proposition~4.1, printed p.~184 of \cite{GGGD2023}, asserts
equivalence of real and complexified invariant-subspace existence for
real Banach lattices whose order is induced by an unconditional basis.
It imposes no positivity or band-diagonal assumption on the operator.
The global convention in its Section~1.1 restricts the spaces to such
lattices; it defines $\mathcal L(X)$ as all bounded linear operators
and introduces positivity separately.  The surrounding Section~4 does
not impose an additional standing positivity assumption.
The paragraph preceding it cites Chapter~9 of the book by Chalendar
and Partington as general background on moment sequences; no proof
of the proposition is supplied there.

The construction above also satisfies that space hypothesis.
On $Y_{\mathbb R}$ the ordered real basis
$(e_1,ie_1,e_2,ie_2,\ldots)$ is $1$-unconditional and its coordinate
order is a Banach lattice order: the norm is
$\sum_n\sqrt{a_n^2+b_n^2}$.  Alternatively use the transported example
on $\ell_1(\mathbb R)$.  It therefore contradicts implication
\textup{(i)$\Rightarrow$(ii)} of that proposition as printed.
The elementary opposite implication, obtained by complexifying a real
invariant subspace, remains valid.  Our example is not a positive
band-diagonal operator, so it does not by itself contradict the positive
band-diagonal existence theorems elsewhere in that paper.

The original formulation \cite[Conjecture~3]{AAST2005} has also been
checked with its preceding assumptions: it concerns arbitrary bounded
operators on infinite-dimensional separable real spaces, with no
additional spectral restriction.  Its Lemma~5 describes a commuting
closed complex structure as the obstruction to the conjecture; here
that structure is the bounded, everywhere-defined $J$ recovered above.
The proof provides the counterexample independently of the conflicting
proposition.  Priority and the reason for the published discrepancy
have not been established in this investigation.


\begin{thebibliography}{99}
\small
\setlength{\itemsep}{3pt}

\bibitem{AAST2005}
Y.~A.~Abramovich, C.~D.~Aliprantis, G.~Sirotkin, and V.~G.~Troitsky,
\emph{Some open problems and conjectures associated with the invariant
subspace problem}, Positivity \textbf{9} (2005), 273--286.
\href{https://doi.org/10.1007/s11117-004-3786-9}{doi:10.1007/s11117-004-3786-9}.
Conjecture 3.

\bibitem{Read1986}
C.~J.~Read,
\emph{A short proof concerning the invariant subspace problem},
J. London Math. Soc. (2) \textbf{34} (1986), 335--348.
\href{https://doi.org/10.1112/jlms/s2-34.2.335}{doi:10.1112/jlms/s2-34.2.335}.

\bibitem{GR2014}
S.~Grivaux and M.~Roginskaya,
\emph{A general approach to Read's type constructions of operators
without non-trivial invariant closed subspaces},
Proc. London Math. Soc. (3) \textbf{109} (2014), 596--652.
\href{https://doi.org/10.1112/plms/pdu012}{doi:10.1112/plms/pdu012}.
\href{https://arxiv.org/abs/1301.6143}{arXiv:1301.6143}, Theorem 1.1.

\bibitem{GGGD2023}
E.~A.~Gallardo-Guti\'errez and F.~J.~Gonz\'alez-Do\~na,
\emph{Band-diagonal operators on Banach lattices: Matrix dynamics and
invariant subspaces}, Linear Algebra Appl. \textbf{674} (2023), 155--191.
\href{https://doi.org/10.1016/j.laa.2023.05.017}{doi:10.1016/j.laa.2023.05.017}.
The conflict discussed in Section~\ref{sec:920} concerns Proposition 4.1.
\end{thebibliography}
\end{document}
